\section{Experiments}

\subsection{Setup}

\paragraph{Datasets and models.}
We evaluate on CIFAR-10, Tiny ImageNet-200,
ImageNet-1K 64x64 HF, strict-source downsampled ImageNet-64, FFHQ-64,
AFHQv2-64, ImageNet-256 10\%, and full ImageNet-256. The HF ImageNet-64 row
and the strict-source downsampled ImageNet-64 row are reported separately; older
HF-derived results are not silently renamed as strict-source results.
All matched experiments use small unconditional models and short training
budgets. The primary model is K8 light \method{} at 64px and K4 at 256px. We
compare it with a parameter-matched direct dense predictor, an endpoint-only
factorized control, explicit objective/prediction/synthesis ablations, Improved
DDPM, and EDM. FlexTok/TiTok reconstruction and official D-AR/MAR/ReTok 50K
results are protocol-separated secondary context.

\paragraph{Protocols.}
The P0 rows match dataset, resolution, split, and nominal optimizer steps;
batch size, nominal images seen, parameter count, and NFE are reported because
the external methods are not compute matched. The direct dense control is the
strictest factorization control: it is parameter matched and uses the same
training budget, but is not claimed to be FLOP or latency matched.
The broad matrix uses locked per-dataset evaluation sizes. Repeated and
ablation claims use seeds 103/139 and 1,024 validation images at fixed $t=500$.
All power-sweep models are evaluated against the same progress-$1.5$ path.

\paragraph{Metrics and evidence tiers.}
We report endpoint PSNR and MSE, path AUC, effective token count, and zero-input
energy. Lower path AUC means better alignment of all prefixes
to the prescribed partial-denoising path; endpoint metrics alone do not measure
that interface. Frechet-style values are internal comparison metrics rather
than SOTA FID claims. The locked evidence matrix contains 80 matched P0 cells,
16 tokenizer eval-only cells, and 24 explicitly protocol-blocked cross-task
cells; full settings and provenance are in the appendix.

\subsection{Results}

\paragraph{All-dataset evidence.}

Table~\ref{tab:cofitok-diagnostics-current} isolates the supported claim. C is
\method{}, D is the parameter-matched direct dense predictor, and E is the
endpoint-only factorized control. \method{} has lower path AUC than E on all
eight datasets and exactly preserves the restricted zero-token contract. Its
endpoint PSNR is close to, but not uniformly better than, direct dense.

\begin{table*}[t]
  \centering
  \scriptsize
  \caption{Diagnostics for the supported \method{} claim. Higher PSNR is better;
  lower path AUC and zero ratio are better.}
  \label{tab:cofitok-diagnostics-current}
  \resizebox{\textwidth}{!}{%
  \begin{tabular}{lrrrrrrrr}
    \toprule
    Dataset & C PSNR & D PSNR & C path & E path & C effK & E effK & C zero & deep-$S_k$ zero \\
    \midrule
    CIFAR-10 & 14.245 & 14.407 & 0.0967 & 0.7906 & 6.84 & 5.24 & 0.0000 & 0.0619 \\
Tiny-IN & 14.348 & 14.453 & 0.0731 & 1.1177 & 6.80 & 5.57 & 0.0000 & 0.0505 \\
ImageNet-64 HF & 14.897 & 15.004 & 0.0672 & 1.1120 & 6.79 & 5.33 & 0.0000 & 0.0522 \\
ImageNet-64 strict & 14.521 & 14.656 & 0.0705 & 1.0502 & 6.82 & 5.36 & 0.0000 & 0.0535 \\
FFHQ-64 & 15.024 & 15.126 & 0.0676 & 1.0830 & 6.82 & 5.41 & 0.0000 & 0.0518 \\
AFHQv2-64 & 15.167 & 15.154 & 0.0653 & 1.0185 & 6.82 & 5.45 & 0.0000 & 0.0561 \\
ImageNet-256 10\% & 13.627 & 13.490 & 0.1416 & 1.5686 & 3.36 & 3.67 & 0.0000 & 0.0179 \\
ImageNet-256 & 13.809 & 13.770 & 0.1325 & 1.5003 & 3.37 & 3.67 & 0.0000 & 0.0185\\
    \bottomrule
  \end{tabular}%
  }
\end{table*}

\begin{figure*}[t]
  \centering
  \begin{minipage}[t]{0.49\textwidth}
    \centering
    \includegraphics[width=\linewidth]{figure2a_tiny_prefix.png}
    \par\smallskip
    {\footnotesize\textbf{(a)} Tiny-IN: endpoint-only (left), \method{} (right).}
  \end{minipage}\hfill
  \begin{minipage}[t]{0.49\textwidth}
    \centering
    \includegraphics[width=\linewidth]{figure2b_imagenet64_prefix.png}
    \par\smallskip
    {\footnotesize\textbf{(b)} ImageNet-64: endpoint-only (left), \method{} (right).}
  \end{minipage}

  \medskip
  \begin{minipage}[t]{0.49\textwidth}
    \centering
    \includegraphics[width=\linewidth]{figure2c_restricted_synthesis.png}
    \par\smallskip
    {\footnotesize\textbf{(c)} Restricted $S_k$: ordered (left), shuffled (right).}
  \end{minipage}\hfill
  \begin{minipage}[t]{0.49\textwidth}
    \centering
    \includegraphics[width=\linewidth]{figure2d_deep_synthesis.png}
    \par\smallskip
    {\footnotesize\textbf{(d)} Deep $S_k$: ordered (left), shuffled (right).}
  \end{minipage}
  \caption{%
  Qualitative prefix and degeneration evidence assembled from four standalone
  subfigure files. In (a--b), each half shows the same examples across the
  prefix path; \method{} produces more organized intermediate denoising than
  endpoint-only factorization. In (c--d), shuffling breaks sample alignment,
  while deep synthesis can carry structure that the restricted, zero-preserving
  operator cannot. Quantitative path and zero-token results are reported in
  Tables~\ref{tab:cofitok-diagnostics-current} and~\ref{tab:component-ablation}.}
  \label{fig:final-visual}
\end{figure*}

\paragraph{Endpoint quality and prefixes.}

At the same 20k training budget, K8 light \method{} nearly matches the
endpoint-only factorized control in endpoint $x_0$ MSE at $t=500$ while greatly
improving path alignment.

\begin{table*}[t]
  \centering
  \caption{Matched 20k two-seed comparison (mean $\pm$ sample standard
  deviation over seeds 103 and 139). Endpoint $x_0$ MSE and path AUC use
  1,024 validation images at $t=500$; lower is better. Higher is better for
  effective token count.}
  \label{tab:fair20k}
  \small
  \resizebox{\textwidth}{!}{%
  \begin{tabular}{llrrrr}
    \toprule
    Dataset & Method & Endpoint $x_0$ MSE $\downarrow$ & Path AUC $\downarrow$ & Eff. K $\uparrow$ & Zero ratio $\downarrow$ \\
    \midrule
  Tiny ImageNet-200 & Endpoint-only factorized & $0.1253 \pm 0.0001$ & $1.0098 \pm 0.1430$ & $5.312 \pm 0.049$ & $0.0000 \pm 0.0000$ \\
Tiny ImageNet-200 & CoFiTok & $0.1296 \pm 0.0015$ & $0.0378 \pm 0.0007$ & $6.766 \pm 0.015$ & $0.0000 \pm 0.0000$ \\
ImageNet-64 HF & Endpoint-only factorized & $0.1096 \pm 0.0006$ & $0.9984 \pm 0.1210$ & $5.169 \pm 0.108$ & $0.0000 \pm 0.0000$ \\
ImageNet-64 HF & CoFiTok & $0.1151 \pm 0.0016$ & $0.0379 \pm 0.0037$ & $6.752 \pm 0.012$ & $0.0000 \pm 0.0000$\\
    \bottomrule
  \end{tabular}
  }
\end{table*}

Across the four paired dataset-seed comparisons, \method{} lowers path AUC in
4/4 cases, with a mean relative reduction of 96.18\%. This structural gain has
a measurable endpoint tradeoff: endpoint $x_0$ MSE increases by 4.20\% on average and
is not lower in any of the four pairs. Figure~\ref{fig:final-visual} shows the
corresponding qualitative prefix grids.
In the full-validation deterministic sweep, K8 light also improves MSE AUC and
low-res Frechet proxy AUC on both Tiny ImageNet-200 and ImageNet-64 HF, while
the endpoint-only factorized control keeps the lower endpoint MSE.

\paragraph{Learned order and ImageNet-256 scaling.}

Order ablations keep the same component sum at the endpoint but change the
prefix trajectory. This makes order evaluation a direct test of whether prefixes
are organized rather than an endpoint-quality test.

Table~\ref{tab:imagenet256-confirmatory-order} gives the preregistered
ImageNet-256 K4 20k two-seed test. Besides a fixed random permutation and
reverse order, it exhaustively evaluates all 24 component permutations over
1,024 images and reports a paired-bootstrap confidence interval for the mean of
the 23 non-identity orders. Because every order sums the same components, the
endpoint is preserved; the test isolates the organization of the prefix path.

\begin{table*}[t]
  \centering
  \scriptsize
  \caption{ImageNet-256 K4 20k confirmatory order and endpoint results over
  1,024 validation images at $t=500$.}
  \label{tab:imagenet256-confirmatory-order}
  \resizebox{\textwidth}{!}{%
  \begin{tabular}{rrrrrrrrr}
    \toprule
    Seed & E path & C ordered & non-ID mean & reverse & rank/24 & $\Delta$ CI low & D endpoint & C endpoint \\
    \midrule
    103 & 1.7879 & 0.0664 & 0.3157 & 0.7192 & 1 & 0.2489 & 0.1234 & 0.1133 \\
139 & 1.4401 & 0.0678 & 0.3082 & 0.6821 & 1 & 0.2400 & 0.1244 & 0.1179\\
    \bottomrule
  \end{tabular}%
  }
\end{table*}

The preregistered confirmatory test passes all seven gates. For both seeds, the
learned order ranks 1st of 24, every non-identity permutation has worse path
AUC, and the paired-bootstrap lower bounds for non-identity minus ordered AUC
are positive (0.2489 and 0.2400). The endpoint component sum is invariant up to
$3\times10^{-15}$ MSE. Relative to the parameter-matched direct dense control,
\method{} also lowers endpoint $x_0$ MSE by 6.73\% on average in this
ImageNet-256 confirmatory setting.

\paragraph{Generation and related-method context.}
Cross-batch validation confirms that endpoint-only factorization remains a
strong endpoint baseline. Generated-sample metrics do not establish stable
unconditional superiority: the matched controls favor endpoint-only on Tiny
ImageNet-200 and ImageNet-64 HF under the internal streamed protocol, and EDM
or direct dense wins the broad Frechet-style rows. Official D-AR/MAR/ReTok 50K
and FlexTok/TiTok reconstruction results remain secondary protocol tiers in the
appendix; none is used as a matched \method{} superiority claim.

\subsection{Ablations}

\paragraph{Matched component protocol.}
We isolate the K8 method components on ImageNet-64 HF with identical 5k-step
training, seeds 103/139, and 1,024-image evaluation at $t=500$. The full model
uses $(\lambda_p,\lambda_c,p)=(0.15,0.30,1.5)$. Endpoint-only removes both path
terms; the next two rows remove one term at a time; simultaneous prediction
removes token feedback; deep synthesis replaces restricted $S_k$ with a biased
nonlinear decoder. Table~\ref{tab:component-ablation} reports mean $\pm$ sample
standard deviation.

\begin{table*}[t]
  \centering
  \small
  \caption{Matched two-seed component ablations on ImageNet-64 HF. Lower is
  better for endpoint MSE, path AUC, and zero ratio; higher effective K indicates
  broader component use. All path values use the same progress-$1.5$ target.}
  \label{tab:component-ablation}
  \resizebox{\textwidth}{!}{%
  \begin{tabular}{lrrrr}
    \toprule
    Variant & Endpoint $x_0$ MSE $\downarrow$ & Path AUC $\downarrow$ & Eff. K $\uparrow$ & Zero ratio $\downarrow$ \\
    \midrule
  CoFiTok (full) & $0.1300 \pm 0.0027$ & $0.0686 \pm 0.0017$ & $6.780 \pm 0.015$ & $0.0000 \pm 0.0000$ \\
Endpoint-only factorized & $0.1275 \pm 0.0019$ & $0.9991 \pm 0.1607$ & $5.366 \pm 0.054$ & $0.0000 \pm 0.0000$ \\
No path-prefix term & $0.1300 \pm 0.0003$ & $0.1269 \pm 0.0096$ & $6.547 \pm 0.027$ & $0.0000 \pm 0.0000$ \\
No path-component term & $0.1302 \pm 0.0012$ & $0.1105 \pm 0.0013$ & $6.752 \pm 0.069$ & $0.0000 \pm 0.0000$ \\
Simultaneous tokens & $0.1314 \pm 0.0016$ & $0.0709 \pm 0.0081$ & $6.829 \pm 0.025$ & $0.0000 \pm 0.0000$ \\
Deep synthesis & $0.1496 \pm 0.0305$ & $0.0463 \pm 0.0051$ & $6.774 \pm 0.032$ & $0.0485 \pm 0.0050$\\
    \bottomrule
  \end{tabular}
  }
\end{table*}

\paragraph{Objective, prediction, and synthesis components.}
Relative to endpoint-only factorization, the full objective lowers mean path AUC
by 93.0\% with a 2.0\% endpoint-MSE cost. Relative to removing only the prefix
or component path term, it lowers path AUC by 45.7\% and 37.9\%, respectively,
while endpoint MSE changes by less than 0.2\%. Both path terms therefore improve
the claimed interface rather than endpoint fit. Sequential feedback is a modest
implementation gain, not a core requirement: full improves the two-seed mean by
2.8\% in path AUC and 1.0\% in endpoint MSE over simultaneous heads, but wins
path AUC on only one seed. Deep synthesis attains lower path AUC yet has zero
ratio $0.0485\pm0.0050$ versus exactly zero for restricted $S_k$; its apparent
gain violates the identifiability contract. Figure~\ref{fig:final-visual}
visualizes zero/shuffle/deep-$S_k$ diagnostics.

\paragraph{Hyperparameter sensitivity.}
Figure~\ref{fig:ablation-hyperparameters} varies the two path-loss weights,
progress power, and token count. Every point is a two-seed mean with sample
standard-deviation bars; the dashed line marks the default. Crucially, the power
sweep is evaluated against one common progress-$1.5$ path rather than allowing
each training target to define its own metric. Raising either path weight
steadily improves path AUC with little endpoint movement; $(0.15,0.30)$ is a
moderate tradeoff rather than the path-only optimum. Power 1.5 is the clear
minimum under the common target. K8 gives the lowest mean path AUC, while K4 has
slightly lower endpoint MSE and K16 adds cost without a 5k-step path gain.

\begin{figure*}[t]
  \centering
  \includegraphics[width=\textwidth]{aaai27_ablation_hyperparameters.pdf}
  \caption{Matched ImageNet-64 HF hyperparameter ablations. Top: denoise-path
  AUC; bottom: endpoint $x_0$ MSE. Curves vary prefix weight $\lambda_p$,
  component weight $\lambda_c$, progress power $p$, and token count $K$.
  Lower is better; bars show sample standard deviation over seeds 103/139.}
  \label{fig:ablation-hyperparameters}
\end{figure*}

\paragraph{Component-separation stress test.}

We treat component decorrelation as a pressure-term ablation rather than a
default objective. The headline \method{} rows use the no-decorrelation light
denoise-path objective because it most directly measures ordered prefix
denoising. On Tiny ImageNet-200, a conservative schedule
(\texttt{weight=0.003}, start step 3000, warmup 1500, full weight at step 4500
of 5000) reduces component correlation on two seeds while keeping the
ordered-prefix penalty small.

Across the two seeds, scheduled decor 0.003 changes ordered path AUC by +0.0036,
mean absolute component cosine by -0.0328, final MSE by -0.0002, and effective
token count by +0.0336. Thus mild separation pressure can make components less
correlated without materially breaking ordered prefix behavior. Stronger decor
weights reduce cosine further but impose larger path-AUC penalties, so we keep
them as stress-test ablations rather than the default \method{} objective.
